\section{Open questions}
\label{sec:6}

We leave three questions for future work.

The first, and the main one this work raises, is the phase diagram above
$\alpha_L$. There the quartic
coefficient of the uniform lattice is negative, so the transition is first
order and the onset $B_c$ is only the spinodal of the normal phase. Locating
the thermodynamic transition and the state it leads to needs the sextic
term on the brane, or the nonlinear condensed branch. In particular,
nothing we have computed bounds the first-order line to $\alpha < \alpha_\star$:
above $\alpha_\star$ the normal phase has no linear charged instability, but
a first-order transition into a condensed phase is not excluded, so
$\alpha_\star$ is established as the end of the instability and not yet as
the end of the vortex lattice.

The second is a proof of triangular minimality at finite coupling.
Numerically the triangular lattice is the global minimum over Bravais
lattices at every coupling we have solved below $\alpha_I$, but
Montgomery's theorem covers only the probe limit: on the brane the weighted
kernel is not completely monotone from $\beta \approx 4$, and below that no
proof is known (section~\ref{sec:5.2}). A proof would need a criterion for such kernels, such as
the sufficient condition of~\cite[theorem~1.1]{Betermin:2016theta}.

The third is the neutral sector. Above $\alpha_\star$ the charged sector
is mode stable: it has no growing quasinormal mode
(section~\ref{sec:4.5}). In this paper the fluctuations of $A^3$ and of
the metric about the brane have been analysed only in the static sector
sourced by the condensate (appendix~\ref{app:response-throat} and section~\ref{sec:5}), and
nonlinear stability is not addressed.
